\section{Widening one: what a condition can see}
\label{sec:logic}

\subsection{The condition language is generated, not designed}

Here is the claim of this section, stated first so the rest can be read as
support for it.

The vocabulary available in a \texttt{where} clause is not a fixed list of
built-in predicates, and nobody sat down and designed it. It is
\emph{constructed from the language definition} that types the channel. Change
the language definition, and the condition language changes with it,
automatically and in a specified way.

That is worth dwelling on because the alternative --- a fixed predicate library
--- is what every other system does, and its failure mode is familiar. A query
language ships with the predicates its designers anticipated. When a developer
needs one that was not anticipated, the escape hatch is a user-defined function,
which is opaque to the optimiser, the index, and the type system. Generated
conditions have no escape hatch because they have no fixed list to escape from.

A generated condition language has three components, and the account of why it
has exactly these three is a realizability argument set out
elsewhere~\cite{Zeus}. Two of them come from the language definition and one is
a choice.

\begin{center}
\small
\begin{tabular}{@{}L{0.22\textwidth}L{0.30\textwidth}L{0.40\textwidth}@{}}
\toprule
\textbf{Component} & \textbf{Where it comes from} & \textbf{What it lets you say} \\
\midrule
structural & one connective per term
constructor in the \texttt{terms} block & ``this term is a send on a channel
satisfying \dots'', ``this splits into two parts sharing nothing'' \\[3pt]
behavioural & one modality per rule in the
\texttt{rewrites} block, and per position it could fire at & ``after one step of
\emph{this} rule, \dots holds'' \\[3pt]
collection & \textbf{a choice}: what
container holds the witnesses & which logical connectives exist at all, and how
they behave \\
\bottomrule
\end{tabular}
\end{center}

\subsection{Structural connectives: pattern matching that can talk about namespaces}

The structural component is the easiest to get a feel for, because it is
recognisably a generalisation of pattern matching.

Every constructor in the \texttt{terms} block yields a connective. If the theory
has \texttt{POutput . n:Name, q:Proc |- n "!" "(" q ")" : Proc}, the logic gets a
way to say ``a send, on a name satisfying $\varphi$, of a term satisfying
$\psi$''. Nested, these compose into descriptions of arbitrary shape.

Two of the generated connectives are worth naming because they do things
ordinary patterns cannot.

\paragraph{Separating conjunction.} Because parallel composition is
associative-commutative --- an equation, from rung two --- the connective it
generates is not ordinary conjunction but a \emph{separating} one: $\varphi \mid
\psi$ holds of a term that splits into two parts, one satisfying $\varphi$ and
one satisfying $\psi$, \emph{with the split being a real partition}. This is the
same idea as separation logic's separating conjunction~\cite{Reynolds02}, and
the same connective the spatial logics for concurrency
carry~\cite{CairesCardelli03}, and it is what lets a
condition say ``there are two of these, and they are not the same one'' --- a
sentence ordinary conjunction cannot express.

\paragraph{The name predicate.} Because names are quoted terms, a condition can
descend from a name into the term it quotes~\cite{NamespaceLogic}: \texttt{@}$\varphi$ holds of a name
whose quoted process satisfies $\varphi$. In developer terms, this means a
condition can talk about \emph{namespaces} --- about the shape of the addresses
involved, not just the payloads. The forager's belief channel and the specimen
channel are distinguishable to a condition by what kind of name they are, which
is how one describes a region of the address space rather than enumerating it.

\subsection{Behavioural connectives: asking about the future}

The behavioural component is the one with no analogue in a query language. For
each rewrite rule, and each position in the term where that rule could fire, the
logic gets a modality $\langle K \rangle \varphi$ --- a refinement of the
modality of Hennessy--Milner logic~\cite{HM85}, indexed by rewrite rule and
position rather than by action~\cite{Hypercube} --- read: \emph{there is a step of
this rule at this position, after which $\varphi$ holds}.

The forager's use for this is immediate. \texttt{Chi(spec)} as written is a
structural question --- does this specimen \emph{look} a certain way. A
behavioural question is different: does this specimen, if poked, \emph{do} a
certain thing. Both are legitimate hypotheses. They are priced very differently.

\begin{remark}[Read or drive]
\label{rem:readordrive}
There are two ways to learn any property of the environment. \emph{Read} it, if
the environment has already computed the property into visible structure --- and
pay only for looking. Or \emph{drive} it, by making the environment take a step
and observing what happens --- and pay for the step. The structural connectives
price the first, the behavioural connectives price the second, and which is
cheaper is a fact about the environment rather than about the observer.
\end{remark}

Remark~\ref{rem:peek} now acquires a sharper form. Driving means acting, and
acting on the thing you are trying to observe changes it. A system that wants
perception to be non-disturbing must arrange that its observational apparatus
uses peeks and its interactions use linear receipts --- and the two must sit in
different namespaces, so that a formula meant to look cannot accidentally claim.
In the forager, reading the specimen's structure and opening the specimen must
be different names, or the assessment eats the subject.

\subsection{The collection component is a dial}

The third component is a choice, and it is the most consequential one in this
section. A logic needs somewhere to put its witnesses --- the things that make a
formula true --- and the container chosen determines which logical operations
even exist.

\begin{center}
\small
\begin{tabular}{@{}L{0.24\textwidth}L{0.28\textwidth}L{0.40\textwidth}@{}}
\toprule
\textbf{Container} & \textbf{Logic you get} & \textbf{What the \texttt{where} can now ask} \\
\midrule
sets (powerset) &
the usual Boolean operations: and, or, not, implication &
``does this hold?'' --- the familiar case, and what the implementation does
today \\[3pt]
quantales~\cite{Rosenthal90,Yetter90} &
the linear connectives~\cite{Girard87} &
``is there \emph{exactly one} of these, to be consumed once?'' --- resource
questions, where using a witness spends it \\[3pt]
join semi-lattices &
coalition logic~\cite{Pauly02,GabbayLosa} &
``does some group of these participants agree?'' --- quorum questions, asked
about a population without running a protocol \\
\bottomrule
\end{tabular}
\end{center}

The distinction between the first two rows is the distinction between counting
and noticing. A Boolean condition, asked whether there are three red items,
answers a question about presence: there \emph{are} red ones. It cannot count
them, because in a set two witnesses of the same kind collapse into one. A linear
condition keeps them apart, because in a quantale they do not collapse. If you
have ever wanted a query that says ``at least three distinct matching rows, and
consuming one of them consumes it for everyone'', the second row is where that
lives.

The third row is stranger and more useful than it sounds. Coalition logic asks
about groups: is there a set of participants such that, together, they can bring
something about. Read as a condition on a running population --- other programs
in the environment --- it asks whether a quorum has formed~\cite{HerdNote}. The forager can use
this to ask whether the specimens around it are of a kind that behaves as a
group, without participating in any protocol they might be running.

\begin{remark}[Same question, three sights]
The reason to present this as a dial rather than three logics is that the
question being asked is the same question, and the theory being asked about is
the same theory. The structural and behavioural connectives are unchanged. What
changes is only what a witness is, and therefore what conjunction and
disjunction do to witnesses. This is why the arrangement is not ad hoc: a
developer is not choosing among logics somebody designed, but among
\emph{observation disciplines} over a structure that was already
there~\cite{ObsDisc}.
\end{remark}

\subsection{What the dial costs}

Turning the dial up is not free, and the price is worth stating in the terms a
developer cares about.

\paragraph{A more expressive condition may not terminate.} The behavioural
modality asks about steps, and a term can step forever. Fixed points over the
modality --- ``it keeps being able to do this'' --- are expressive and
unbounded. The mitigation is that evaluating a guard is metered like everything
else: a budget bounds the check, and an unfunded check is a non-event by
\S\ref{sec:cost}, so it fails safe.

\paragraph{Restricting to make checking terminate loses something real.} A
developer who wants guards that always answer can restrict to a fragment ---
fewer connectives, or a smaller slice of the theory. That works. What is
surrendered is \emph{adequacy}: the property that the logic can distinguish any
two programs that actually behave differently. In general, a logic whose model
checking always terminates cannot be adequate for a Turing-complete theory. This
is not a defect of the construction; it is the shape of the trade, and the dial
exists precisely so it can be made deliberately rather than by accident.

\paragraph{Adequacy is a property of the target, not of the generator.} It is
tempting to assume that a machine-generated logic comes with a guarantee. It
does not, and it should not: whether the generated logic is adequate depends on
which connectives were asked for. Without something like a negation and a
conjunction, it will not be. A generated \emph{type system}, for instance, is
generated by the same family of algorithms and is not expected to be adequate at
all~\cite{Hypercube,GSLTintro}.

\subsection{The forager's ladder}

Return to the running example. \texttt{Chi(spec)} was a placeholder for one
hypothesis. The dial says there is a whole ladder of possible hypotheses, from
cheap structural ones to expensive behavioural ones, and the forager can afford
to sit anywhere on it.

That is not a metaphor: it has been measured. In a companion
study~\cite{GameNote}, a learner of
exactly this shape was equipped with a ladder of hypotheses expressed as
\texttt{where} clauses, each priced by the depth of the formula it installs, and
run to exhaustion against a fixed environment.

\begin{center}
\small
\begin{tabular}{@{}L{0.52\textwidth}L{0.40\textwidth}@{}}
\toprule
\textbf{Measurement} & \textbf{What it says} \\
\midrule
The non-modal rungs delivered $92.3\%$ of the total gain for $34.3\%$ of the
total cost &
structural sight is the bargain; behavioural sight is the luxury \\[3pt]
The single best purchase on the whole ladder was a \emph{name-level} predicate
at modal depth zero &
namespaces are cheap to look at and highly informative \\[3pt]
The rung that quantified over all replies to all moves was the worst purchase
measured &
alternating quantification over the future is where cost explodes \\[3pt]
A learner holding the \emph{complete, correct} theory was dominated by a
learner holding a cheaper wrong one &
truth buys safety, not yield; and the environment, not the learner, decides
whether truth is affordable \\
\bottomrule
\end{tabular}
\end{center}

The last row is the one to carry forward. It says that the right position on the
dial is not a property of the problem, and certainly not a matter of rigour. It
is an economic question about the environment the program is deployed into, and
it is answerable --- because both the sight and the acting are priced in the same
currency.
